\subsection{Système de séparation}
\label{subsec:sys_sepa}

\subsubsection{Objectifs mécaniques et intégration structurelle}

L'interface de séparation entre le premier et le second étage a un double rôle :
\begin{itemize}
    \item assurer un verrouillage mécanique fiable des deux étages pendant la propulsion du premier étage, sans jeu perceptible ;
    \item limiter la flèche globale du lanceur en se comportant au plus proche d'un tronçon de tube continu, malgré l'intégration du
    mécanisme de séparation.\\
\end{itemize}

Pour répondre à ces objectifs, le mécanisme de séparation est entièrement encapsulé dans un ensemble d'enveloppes structurelles
cylindriques. Les deux enveloppes, solidaires respectivement du premier et du second étage, reprennent les efforts de compression et de
flexion tandis que le mécanisme interne ne reprend que les efforts de verrouillage en traction. La continuité de section et l'utilisation
d'un alliage d'aluminium 7075 garantissent une bonne rigidité en flexion, ce qui limite la flèche globale du lanceur et les charges
parasites sur les autres sous-systèmes.\\

L'enveloppe inférieure intègre en outre deux ouvertures diamétralement opposées, usinées dans la jupe du premier étage. Ces évents
permettent d'évacuer les gaz en cas d'allumage intempestif du moteur du second étage et de satisfaire à la règle CP7 du cahier des charges
Fusex bi-étage \cite{CDC_BiEtage}. Ils sont dimensionnés de manière à éviter toute mise sous pression de l'interface inter-étages et à
annuler la poussée résultante en orientant les jets de gaz de façon symétrique.

\subsubsection{Architecture mécanique générale}

Le mécanisme est actionné par un servomoteur Feetech, identique à celui utilisé pour les trappes parachute, afin de mutualiser les
références et simplifier la maintenance. Le servo est accouplé à une vis sans fin qui entraîne un ensemble bielle-loquet, puis, en fin de
course, un système de libération de ressorts de séparation servant de secours et de moyen de tests au sol.\\

L'interface inter-étages est constituée de :
\begin{itemize}
    \item une pièce structurelle supérieure, solidaire du second étage, usinée en aluminium 7075 et intégrant les logements des quatre
    tenons d'accrochage ;
    \item une pièce structurelle inférieure, solidaire du premier étage, qui reçoit le mécanisme interne (vis sans fin, bielles et loquets),
    les ressorts de séparation (backup), la caméra Arducam ainsi que les ouvertures d'évent pour l'évacuation des gaz.\\
\end{itemize}

Le coeur du mécanisme est une vis sans fin acier qui engrène sur un écrou en laiton. Ce couple vis sans fin / écrou offre un effet
autobloquant qui empêche tout retour de mouvement imposé par les vibrations, par les ressorts ou par les efforts aérodynamiques. Cet écrou
entraîne un chariot guidé linéairement, sur lequel sont articulées quatre paires de bielles réparties à 90 degrés. Chaque paire de bielles
pilote un loquet radial qui vient s'engager dans une fenêtre de la pièce supérieure.\\

A l'état verrouillé, les quatre loquets travaillent en cisaillement et assurent la reprise des efforts axiaux entre les étages, tandis que
les pièces cylindriques reprennent principalement les efforts de flexion.

\subsubsection{Cinématique de déverrouillage en vol}

La séquence nominale de déverrouillage est la suivante :
\begin{enumerate}
    \item le séquenceur du premier étage commande le servomoteur Feetech après la fin de propulsion du premier étage, lorsque les aérofreins
    sont prêts à être déployés ;
    \item le servomoteur fait tourner la vis sans fin, ce qui déplace l'écrou et le chariot dans une translation axiale contrôlée ;
    \item le mouvement du chariot tire les quatre bielles, qui font pivoter les loquets vers l'intérieur ;
    \item lorsque la butée de fin de course de déverrouillage est atteinte, les loquets sont complètement dégagés des fenêtres de la bague
    supérieure.\\
\end{enumerate}

A partir de cet instant, les deux étages ne sont plus liés que par le contact géométrique des surfaces de centrage. La différence de traînée
induite par le déploiement des aérofreins du premier étage génère alors spontanément la séparation des deux étages, sans choc ni poussée
supplémentaire. La nature autobloquante de la vis sans fin interdit tout reverrouillage involontaire et garantit une position mécanique bien
définie, même en cas de coupure d'alimentation du servomoteur.

\subsubsection{Séparation mécanique de secours par ressorts}

Pour satisfaire l'exigence de séparation au sol et garantir une dissociation des étages même en cas de défaillance du système
d'aérofreinage, le mécanisme intègre une seconde fonction de séparation active par ressorts.\\

Au-delà de la butée de déverrouillage, la poursuite de la rotation de la vis sans fin entraîne le chariot sur une deuxième portion de
course. Celui-ci vient alors appuyer sur deux pièces intermédiaires précontraintes par de petits ressorts de retenue. Lorsque l'effort
transmis par le chariot dépasse la précontrainte, ces pièces se dégagent et libèrent deux ressorts principaux de séparation, disposés
axialement dans le premier étage.\\

Dans leur état armé, ces ressorts sont comprimés et assurent un effort nul sur les étages grâce au verrouillage mécanique des pièces
intermédiaires. Une fois libérés, ils exercent une force axiale qui repousse les deux étages l'un de l'autre, produisant :
\begin{itemize}
    \item une séparation franche et reproductible lors des essais de séparation au sol ;
    \item une séparation de secours en vol si l'aérofreinage ne génère pas une divergence suffisante des trajectoires.\\
\end{itemize}

Le séquenceur différencie explicitement ces deux modes :
\begin{itemize}
    \item en mode nominal, la consigne de rotation du servomoteur ne parcourt que la course de déverrouillage des loquets ;
    \item en mode test au sol ou secours, la consigne autorise la course complète intégrant la libération des ressorts.
\end{itemize}

\subsubsection{Gestion d'un allumage intempestif du moteur du second étage}

Conformément à la règle CP7 du cahier des charges Fusex bi-étage, l'interface inter-étages doit rester ouverte en cas d'allumage intempestif
du moteur du second étage, de manière à éviter la mise sous pression du volume inter-étage et à ne pas créer de poussée dissymétrique au sol
ou en vol.\\

Pour satisfaire cette exigence, deux ouvertures rectangulaires sont usinées dans la jupe inférieure, de part et d'autre de l'axe. Elles sont
disposées à 180 degrés, de façon à orienter les jets de gaz symétriquement et à annuler la poussée résultante. La surface libre totale des
évents est choisie supérieure à la section de sortie du moteur du second étage, ce qui limite la surpression maximale dans l'inter-étage à
une valeur compatible avec la tenue mécanique des pièces structurelles.\\

Les bords des ouvertures sont rayonnés afin de réduire les concentrations de contraintes mécaniques et thermiques et d'éviter toute arête
tranchante susceptible d'endommager le câblage de la ligne de mise à feu. La géométrie interne de l'interface est par ailleurs conçue sans
volume fermé ni cloisonnement susceptible de piéger les gaz.

\subsubsection{Matériaux et procédés de fabrication}

Les pièces structurelles principales (pièce supérieure, pièce inférieure, noyau central) sont usinées dans un alliage d'aluminium 7075 de
grade aéronautique, choisi pour son rapport résistance mécanique / masse et sa bonne tenue en fatigue. Ces pièces sont confiées à un
sous-traitant d'usinage \acrshort{CNC} 5 axes (JLCCNC), après comparaison de plusieurs fournisseurs sur les critères suivants :
\begin{itemize}
    \item précision d'usinage compatible avec les jeux fonctionnels du mécanisme ;
    \item coût de production ;
    \item délais de fabrication compatibles avec le planning des \acrshort{RCE}.\\
\end{itemize}

Les composants de type bielles, loquets et certaines brides sont réalisés en tôlerie d'aluminium, découpés et usinés sur la \acrshort{CNC} du
\gls{Skylab}. Cette approche permet de :
\begin{itemize}
    \item valider rapidement plusieurs itérations de géométrie ;
    \item maîtriser les coûts en réalisant en interne les pièces les plus simples ;
    \item ajuster facilement les épaisseurs et rayons pour optimiser la masse.\\
\end{itemize}

Enfin, certains éléments non structuraux (guides de câbles, caches, supports de connecteurs) sont imprimés en PLA renforcé fibre de carbone.
Ces pièces assurent uniquement des fonctions de maintien ou de guidage et ne participent pas à la reprise d'efforts principaux.

\subsubsection{Instrumentation vidéo de la séparation}

Le bloc de séparation intègre également une instrumentation vidéo dédiée à l'observation de la séparation en vol. Une caméra Arducam 12
mégapixels est montée au centre du noyau structurel, dans la partie solidaire du premier étage. Son champ de vision est orienté vers le
haut, exactement dans l'axe du propulseur du second étage.\\

Ce positionnement permet :
\begin{itemize}
    \item d'observer la divergence des deux étages ;
    \item de visualiser l'allumage du second étage et la mise en poussée ;
    \item de disposer d'un retour vidéo précieux pour l'analyse post-vol et la validation de la modélisation de la séparation.\\
\end{itemize}

Les données de la caméra sont enregistrées localement pendant le vol.

\subsubsection{Illustrations du mécanisme de séparation}

La figure \ref{fig:separation_coupe_ouverte} présente une vue en coupe longitudinale du mécanisme en position de déverrouillage, tandis que
la figure \ref{fig:separation_coupe_fermee} montre la configuration verrouillée avec les loquets engagés et les ressorts armés.

% \begin{figure}[h!]
%     \centering
%     \includegraphics[width=0.7\linewidth]{fig/separation_coupe_ouverte.png}
%     \caption{Vue en coupe du mécanisme de séparation inter-étages, position de déverrouillage, loquets dégagés.}
%     \label{fig:separation_coupe_ouverte}
% \end{figure}

% \begin{figure}[h!]
%     \centering
%     \includegraphics[width=0.7\linewidth]{fig/separation_coupe_verrouillee.png}
%     \caption{Vue en coupe du mécanisme de séparation inter-étages, position verrouillée, loquets engagés et ressorts armés.}
%     \label{fig:separation_coupe_fermee}
% \end{figure}

La figure \ref{fig:separation_ensemble} illustre l'intégration de l'ensemble de séparation dans les jupes des étages, avec les enveloppes
structurelles en aluminium 7075, la position de la caméra Arducam et les ouvertures d'évent. Enfin, la figure
\ref{fig:separation_mecanisme_nu} montre le mécanisme nu (bielles, loquets, vis sans fin) pour une meilleure lisibilité de la cinématique.

% \begin{figure}[h!]
%     \centering
%     \includegraphics[width=0.7\linewidth]{fig/separation_ensemble_explose.png}
%     \caption{Vue d'ensemble éclatée du bloc de séparation et des jupes d'étage, montrant l'intégration structurelle et les évents
%     d'évacuation des gaz.}
%     \label{fig:separation_ensemble}
% \end{figure}

% \begin{figure}[h!]
%     \centering
%     \includegraphics[width=0.55\linewidth]{fig/separation_mecanisme_nu.png}
%     \caption{Mécanisme de séparation nu avec les quatre bielles, les loquets et la vis sans fin actionnée par le servomoteur Feetech.}
%     \label{fig:separation_mecanisme_nu}
% \end{figure}

\newpage